\documentclass[11pt]{article}
\usepackage[a4paper,margin=1in]{geometry}
\usepackage{booktabs}
\usepackage{array}
\usepackage{amsmath,amssymb}
\usepackage{xcolor}
\usepackage{fancyhdr}
\usepackage{enumitem}
\usepackage[colorlinks=true,linkcolor=black,urlcolor=black,citecolor=black]{hyperref}

\definecolor{secblue}{RGB}{31,73,125}
\definecolor{boxgray}{RGB}{244,246,249}
\usepackage{sectsty}
\sectionfont{\color{secblue}\large}
\subsectionfont{\color{secblue}\normalsize}
\usepackage[most]{tcolorbox}

\pagestyle{fancy}\fancyhf{}
\renewcommand{\headrulewidth}{0.4pt}
\lhead{\small Bruise Segmentation}\rhead{\small Project Report}\cfoot{\thepage}

\newcolumntype{R}{>{\raggedleft\arraybackslash}p{1.9cm}}
\setlength{\parindent}{0pt}\setlength{\parskip}{5pt}
\newtcolorbox{keybox}{colback=boxgray,colframe=secblue,boxrule=0.6pt,arc=2pt,
  left=8pt,right=8pt,top=6pt,bottom=6pt}

\begin{document}

\begin{center}
  {\LARGE\bfseries Fairness-Aware Bruise Segmentation}\\[3pt]
  {\large A Project Report: Core Models, Baselines, Capacity Scaling, and Knowledge Distillation}\\[3pt]
  {\small White-light bruise photographs \, $\cdot$ \, SegFormer / YOLO / U-Net / DeepLab \, $\cdot$ \, ITA skin-tone groups}
\end{center}

\vspace{2pt}

\begin{keybox}
\textbf{Scope.} This report documents the full experimental arc of the project: (1) the five
core models trained on Colab (SegFormer-B2 teacher, SegFormer-B0 direct/distilled, YOLO26n-sem
direct/distilled); (2) two classical CNN baselines (U-Net, DeepLabV3+); (3) the largest-capacity
SegFormer-B5; and (4) the ongoing knowledge-distillation study. All models use one fixed protocol
(697/134/185 subject-grouped split, $640{\times}640$, scored against a 2-of-3 consensus mask).
Numbers are reported honestly with multi-seed variability where available.
\end{keybox}

\section{Task and Protocol}

The goal is pixel-level segmentation of bruises in white-light (WL) photographs, evaluated for
\emph{equity across skin tone}. The dataset is a multi-labeler study: each image was annotated by
three raters, and the segmentation target is their \textbf{2-of-3 majority-vote consensus mask}.

\begin{center}\small
\begin{tabular}{ll}
\toprule
Split & 697 train / 134 val / 185 test images (subject-grouped, no subject leakage) \\
Test subjects & 28 \\
Resolution & $640\times640$ (bilinear image, nearest mask) \\
Skin tone & Individual Typology Angle (ITA), 5 groups (Light $\to$ Dark) \\
Selection & best epoch on val Dice; decision threshold swept on val; scored \emph{once} on test \\
Primary metrics & median/mean Dice, \textbf{complete-miss rate} (blank mask on a real bruise) \\
\bottomrule
\end{tabular}
\end{center}

\textbf{Inter-annotator context.} On the same 185 test images, pairwise annotator agreement (Dice)
ranges $0.58$--$0.75$ (three-pair mean $\approx0.64$); annotator-vs-consensus ranges $0.70$--$0.87$.
These indicate substantial human labeling variability and are reported as \emph{context} only ---
annotator-vs-annotator and model-vs-consensus Dice are not directly comparable.

\section{Phase 1 --- Core Models (Colab)}

Five models trained with one shared recipe (SegFormer: Dice+BCE, AdamW backbone/head LR split
$6\!\times\!10^{-5}/6\!\times\!10^{-4}$, warmup$\to$poly, AMP, effective batch 8; YOLO: native
Ultralytics recipe). Distillation used calibrated response-based KD (temperature-scaled soft
targets). Results are the \textbf{mean over 3 seeds}; YOLO is reported on its native-argmax path.

\begin{center}\small
\begin{tabular}{l R R R R}
\toprule
Model & Params (M) & Mean Dice & Median Dice & Miss (\%) \\
\midrule
SegFormer-B2 teacher       & 27.35 & $0.7651$ & $0.8113$ & $0.00$ \\
SegFormer-B0 direct        & 3.71  & $0.7600$ & $0.8113$ & $0.54$ \\
SegFormer-B0 distilled     & 3.71  & $0.7639$ & $0.8114$ & $0.18$ \\
YOLO26n-sem direct         & 1.63  & $0.6911$ & $0.8021$ & $7.57$ \\
YOLO26n-sem distilled      & 1.63  & $0.6676$ & $0.7658$ & $7.57$ \\
\bottomrule
\end{tabular}
\end{center}

\noindent\textbf{Findings.}
\begin{itemize}[leftmargin=1.4em,itemsep=1pt]
  \item \textbf{Distillation preserves/matches the teacher at $7.4\times$ fewer parameters}: the
  3.7M distilled B0 (mean $0.7639$) is statistically indistinguishable from the 27M B2 teacher
  (mean $0.7651$), and beats the direct-trained B0 --- the compact student loses essentially nothing.
  \item \textbf{Complete-miss is the discriminating metric, and it separates the architectures}:
  all SegFormers miss $\le0.5\%$ of bruises entirely, whereas YOLO misses $\approx7.6\%$ --- a
  high-consequence failure for an injury-documentation tool. YOLO's higher \emph{median} Dice is
  misleading given this miss rate.
  \item YOLO's ImageNet-normalization/scale handling was a known pitfall; native-argmax is the sole
  correct reporting path (the custom-/255 path is excluded).
\end{itemize}

\section{Phase 2 --- Classical CNN Baselines}

To place the transformers against standard segmentation CNNs, U-Net and DeepLabV3+ (both
ResNet-50 encoders, \texttt{segmentation\_models\_pytorch}) were trained with the same split,
resolution, and evaluation (3 seeds).

\begin{center}\small
\begin{tabular}{l R R R R}
\toprule
Model & Params (M) & Mean Dice & Median Dice & Miss (\%) \\
\midrule
U-Net (ResNet-50)       & $\sim$32 & $0.7526$ & $0.8260$ & $3.24$ \\
DeepLabV3+ (ResNet-50)  & $\sim$40 & $0.7453$ & $0.8140$ & $1.80$ \\
\bottomrule
\end{tabular}
\end{center}

\textbf{Finding.} The CNN baselines are competitive --- U-Net's \emph{median} Dice ($0.8260$) is
actually the highest of the pre-distillation models --- but they miss more bruises entirely
($1.8$--$3.2\%$) than the SegFormers ($\le0.5\%$). No single architecture dominates on all metrics.

\section{Phase 3 --- Capacity Scaling: SegFormer-B5}

SegFormer-B5 ($\sim$85M, the largest MiT encoder) was trained with the identical recipe to test
whether more capacity raises the ceiling. (Result below is a single seed; the 3-seed run is
in progress.)

\begin{center}\small
\begin{tabular}{l R R R R}
\toprule
Model & Params (M) & Mean Dice & Median Dice & Miss (\%) \\
\midrule
SegFormer-B5 (1 seed) & $\sim$85 & $0.7829$ & $0.8276$ & $0.00$ \\
\bottomrule
\end{tabular}
\end{center}

\begin{keybox}
\textbf{Central observation across Phases 1--3.} Across seven very different architectures
(1.6M YOLO $\to$ 85M B5; CNN and transformer), median Dice sits in a narrow band ($0.766$--$0.828$)
and B5 ($0.828$) barely exceeds the 3.7M distilled B0 ($0.811$) and the U-Net baseline ($0.826$).
\emph{Increasing model capacity alone gives diminishing returns under the current data and
evaluation protocol} --- the bottleneck appears more data/label-related than capacity-related. This
is stated as an observation, not a proven ``ceiling'': the validation oracle (\S5) shows there is
still headroom that a \emph{combination} of models can reach.
\end{keybox}

\section{Phase 4 --- Knowledge-Distillation Study (current)}

Given the above, the distillation study is framed not around beating the teacher on Dice (which is
flat) but around: \textbf{can a compact 3.7M student preserve overall quality while reducing
high-consequence failures and deployment cost?} It uses the best-of-3-seed B5 (selected on val,
temperature $T{=}2.20$) and the B2 teacher, distilling into B0.

\subsection{The multi-teacher gate (validation oracle)}
Before committing to multi-teacher methods, a subject-cluster bootstrap on \emph{validation} tested
whether B2 and B5 are complementary:

\begin{center}\small
\begin{tabular}{l R}
\toprule
Quantity (validation) & Value \\
\midrule
Spearman(B2 Dice, B5 Dice) & $0.82$ \\
Images B5 beats B2 by $>0.10$ / vice-versa & $18$ / $12$ \\
Oracle (best-teacher-per-image) gain over best single & $+0.0258$ \\
\quad 95\% CI & $[0.0165,\ 0.0362]$ \\
\quad $P(\text{gain}>0)$ & $1.00$ \\
Adaptive gate & \textbf{PASS} (run multi-teacher) \\
Reliability weight $\mathrm{rel}_b$ (B5) & $0.634$ \\
\bottomrule
\end{tabular}
\end{center}

\textbf{Interpretation.} B2 and B5 make genuinely different errors; a perfect per-image selector
gains $\sim$2.6 Dice points over either single teacher (CI clear of zero). So the single-model
Dice is \emph{not} a hard ceiling --- combining complementary teachers has real, measured headroom.
This decision was made on validation only; the test set is untouched.

\subsection{Methods under test}
All KD terms are additions to the calibrated response loss
$\mathcal{L}=\alpha\,\mathrm{DiceBCE}(z_s,y)+(1-\alpha)\,\mathrm{BCE}(z_s,\sigma(z_t/T))$.
The mix $\alpha$ is searched on validation (grid; Optuna if available).

\begin{center}\small
\begin{tabular}{l l}
\toprule
Arm & What it transfers \\
\midrule
Response KD & calibrated teacher probabilities (baseline) \\
CWD (ICCV'21) & per-channel spatial distributions (canonical seg-KD baseline) \\
BPKD (WACV'24) & boundary vs body knowledge, separately \\
Angular (WACV'24) & feature \emph{direction} (scale-robust for the B5$\to$B0 gap) \\
DKD (CVPR'22) & decoupled target / non-target logits (lower priority in binary seg) \\
Uniform ensemble & $\tfrac12(p_{B2}+p_{B5})$ \\
\textbf{Reliability-gated (novel)} & per-pixel confidence$\times$disagreement gate, val-derived $\mathrm{rel}_b$ \\
\quad + boundary / + hard / + group / + full & one component added at a time (ablation) \\
\bottomrule
\end{tabular}
\end{center}

The novel method combines the reliability gate with paired WL/ALS potential, failure-risk
weighting, and boundary knowledge --- teacher mixing itself is not new (cf.\ HeteroAKD, AAAI'25),
the combination is.

\subsection{Success hierarchy (pre-registered)}
Because Dice is flat, success is judged in priority order:
(1) \textbf{non-inferiority in subject-level Dice} (margin 0.03);
(2) reduction in \textbf{complete-miss and the low-recall tail} (\%recall$<$0.10, \%Dice$<$0.20,
5th-percentile subject Dice);
(3) improved \textbf{small-bruise recall};
(4) improved \textbf{worst-ITA-group recall} (exploratory);
(5) lower \textbf{deployment cost}. Every comparison uses the paired subject-cluster bootstrap
(28 subjects) against the B2$\to$B0 reference. TTA is reported separately as an inference-only
baseline.

\subsection{Early result}
The first completed arm, B5$\to$B0 response KD (validation-searched $\alpha=0.625$):

\begin{center}\small
\begin{tabular}{l R R R R R}
\toprule
Arm & Mean Dice & Median Dice & Miss (\%) & \%rec$<$.10 & worst-grp Dice \\
\midrule
B5$\to$B0 response & $0.7683$ & $0.8061$ & $0.00$ & $2.2$ & $0.7587$ \\
\bottomrule
\end{tabular}
\end{center}

\noindent In line with expectation: overall Dice tracks the B0-distilled reference, with $0\%$
complete miss --- the interesting comparisons (uniform vs reliability-gated ensemble, and the
failure-tail metrics) come from the arms still training.

\section{Status and Reproducibility}

\textbf{Running now (ORC, one GPU, resumable):} the full arm matrix --- response, CWD, BPKD,
uniform ensemble, reliability-gated adaptive + four ablations, angular, DKD --- each with
epoch-level resume (an interrupted run is never scored or marked done) and validation-based
$\alpha$ search. A single aggregate report then labels each arm
\textbf{REPORTABLE WIN / NON-INFERIOR / INFERIOR} vs the B2$\to$B0 reference on the hierarchy above.

\textbf{Deferred / next:} YOLO distillation from the winning teacher configuration (YOLO is the
real weakness at $\sim$7.6\% miss); a data-merge $2\times2$ factor; and, to actually test the
``are we near the data limit?'' question, data-scaling and re-annotation experiments (distillation
alone cannot answer it).

\textbf{One-line summary of the whole project.}
\begin{keybox}
Seven architectures (CNN and transformer, 1.6M--85M) cluster tightly on Dice, a 3.7M distilled
student already matches a 27M teacher, and validation shows two teachers are complementary with
$\sim$2.6-point oracle headroom --- so the distillation study asks how to transfer that complementary
knowledge into the compact student while cutting the high-consequence failures (complete miss,
worst-group, small-bruise) that Dice hides.
\end{keybox}

\vfill
\begin{center}\footnotesize\itshape
Numbers: Colab core models and CNN baselines are 3-seed means (native-argmax for YOLO);
SegFormer-B5 is a single seed (3-seed in progress). All scored on the 185-image consensus test set
under one fixed protocol. Distillation arms use a locked re-evaluation for cross-arm comparability.
\end{center}

\end{document}
